\chapter{Traçabilité et reproduction des résultats}
\label{app:evidence}

\section{Révisions de référence}

\begin{itemize}
  \item base amont : \GitCommit{5fc0681} ;
  \item révision PCB vérifiée au DRC : \GitCommit{b58eda6} ;
  \item état local et distant incluant la base du HIL : \GitCommit{77a253e} ;
  \item scénario HIL bidirectionnel : version de travail \texttt{node\_node} déployée sur le banc le 16 août 2026.
\end{itemize}

\section{Commandes principales}

\begin{lstlisting}[language=bash,caption={Simulateur de mission : installation, exécution et suite de tests.}]
cd simulator
python -m pip install -r requirements.txt
streamlit run app.py          # tableau de bord sur http://localhost:8501
python -m pytest tests -q     # suite de tests unitaires
\end{lstlisting}

\begin{lstlisting}[language=bash,caption={Contrôle DRC de la carte, exporté en JSON.}]
kicad-cli pcb drc --format json --units mm --severity-error \
  --severity-warning --output caiman_drc.json <chemin>/caiman.kicad_pcb
\end{lstlisting}

\begin{lstlisting}[language=bash,caption={Banc HIL : compilation native, flash de l'ESP32 et exécution bidirectionnelle.}]
cd /home/rasp/Caiman/protocol_hil
cmake -S . -B build && cmake --build build --parallel 4
ctest --test-dir build --output-on-failure

source /home/rasp/esp-idf-v6.0.2/export.sh
cd scenarios/node_node/esp32 && idf.py build
idf.py -p /dev/ttyUSB0 flash
sudo nmcli connection up caiman-hil ifname wlan0
cd /home/rasp/Caiman/protocol_hil
./build/caiman_node_r2 192.168.4.1 4210 --count 5
\end{lstlisting}

Le journal synthétique de l'essai est conservé dans
\path{rapport_stage/evidence/hil_2026-08-16.log}. Les deux tests natifs ont
réussi et le programme HIL s'est terminé avec un code de retour nul après cinq
cycles bidirectionnels.

\section{Journaux complets du banc HIL}

Les deux journaux ci-dessous ont été capturés simultanément le 16 août 2026 lors de la seconde exécution du banc. Le \cref{tab:hil-cycles} en donne la synthèse.

\begin{lstlisting}[
  caption={Sortie complète de \texttt{caiman\_node\_r2 192.168.4.1 4210 -{}-count 5} sur le Raspberry Pi~4.},
  label={lst:hil-r2},
  basicstyle=\ttfamily\scriptsize, frame=single, breaklines=true]
R2 node=2 listening on UDP 4210; peer R1=192.168.4.1 node=1

R2 RX PEER_STATE seq=102762 R1->R2 | x=123.9m y=456.9m depth=12.10m
   bottom=18.55m battery=79.5% link=0% leak=no GNSS=no
R2 FRAME 32B encrypted=41010201916a00083a07039bff1df48a...
R2 TX PEER_STATE seq=32769 R2->R1 x=800.0 y=200.0 depth=10.00
   battery=95.0 link=100%
R2 TX PEER_STATE seq=32769 -> TX_DS

R2 RX PEER_STATE seq=102763 R1->R2 | x=124.4m y=457.1m depth=12.20m
   bottom=18.60m battery=79.0% link=53% leak=no GNSS=no
R2 TX PEER_STATE seq=32770 R2->R1 x=799.7 y=200.4 depth=10.08
   battery=94.5 link=75%
R2 TX PEER_STATE seq=32770 -> TX_DS

R2 RX PEER_STATE seq=102764 R1->R2 | x=124.9m y=457.3m depth=12.30m
   bottom=18.65m battery=78.5% link=100% leak=no GNSS=no
R2 TX PEER_STATE seq=32771 R2->R1 x=799.4 y=200.8 depth=10.16
   battery=94.0 link=75%
R2 TX PEER_STATE seq=32771 -> TX_DS

R2 RX PEER_STATE seq=102765 R1->R2 | x=125.4m y=457.5m depth=12.40m
   bottom=18.70m battery=78.0% link=100% leak=no GNSS=no
R2 TX PEER_STATE seq=32772 R2->R1 x=799.1 y=201.2 depth=10.24
   battery=93.5 link=75%
R2 TX PEER_STATE seq=32772 -> TX_DS

R2 RX PEER_STATE seq=102766 R1->R2 | x=125.9m y=457.7m depth=12.50m
   bottom=18.50m battery=77.5% link=100% leak=no GNSS=no
R2 TX PEER_STATE seq=32773 R2->R1 x=798.8 y=201.6 depth=10.32
   battery=93.0 link=75%
R2 TX PEER_STATE seq=32773 -> TX_DS
\end{lstlisting}

\begin{lstlisting}[
  caption={Événements de diagnostic de l'ESP32 (R1), capturés par \texttt{caiman\_esp\_log 4211}.},
  label={lst:hil-r1},
  basicstyle=\ttfamily\scriptsize, frame=single, breaklines=true]
ESP32 HIL diagnostic events on UDP 4211
R1 TX PEER_STATE seq=102762 payload=8B frame=32B -> TX_DS
R1 RX PEER_STATE seq=32769 R2->R1 x=800.0 y=200.0 depth=10.00
   battery=95.0 link=100%
R1 TX PEER_STATE seq=102763 payload=8B frame=32B -> TX_DS
R1 RX PEER_STATE seq=32770 R2->R1 x=799.7 y=200.4 depth=10.08
   battery=94.5 link=73%
R1 TX PEER_STATE seq=102764 payload=8B frame=32B -> TX_DS
R1 RX PEER_STATE seq=32771 R2->R1 x=799.4 y=200.8 depth=10.16
   battery=94.0 link=73%
R1 TX PEER_STATE seq=102765 payload=8B frame=32B -> TX_DS
R1 RX PEER_STATE seq=32772 R2->R1 x=799.1 y=201.2 depth=10.24
   battery=93.5 link=73%
R1 TX PEER_STATE seq=102766 payload=8B frame=32B -> TX_DS
R1 RX PEER_STATE seq=32773 R2->R1 x=798.8 y=201.6 depth=10.32
   battery=93.0 link=73%
R1 TX PEER_STATE seq=102767 frame=32B -> MAX_RT
\end{lstlisting}

\section{Schéma électronique complet}

Le \cref{fig:schematic-annexe} reproduit le schéma KiCad complet dans une taille exploitable. Les extraits fonctionnels discutés dans le corps du rapport en sont issus.

\begin{figure}[htbp]
  \centering
  \includegraphics[width=\textwidth,height=0.82\textheight,keepaspectratio]{fabrication/caiman_schematic.pdf}
  \caption{Schéma électronique complet de la carte Caiman (microcontrôleur, alimentations, bus d'interfaçage et capteurs).}
  \label{fig:schematic-annexe}
\end{figure}

\section{Preuves complémentaires attendues}
\label{app:preuves}

\begin{itemize}\setlength{\itemsep}{0pt}
  \item journal ou tracé des lectures du banc STM32/ICM-20948 ;
  \item capture datée de \url{https://caimansim.fr} après remise en ligne ;
  \item photographies et relevés de continuité de la carte après réception ;
  \item résultats du bring-up, postérieurs à la soutenance si nécessaire.
\end{itemize}
